% Options for packages loaded elsewhere
\PassOptionsToPackage{unicode}{hyperref}
\PassOptionsToPackage{hyphens}{url}
\PassOptionsToPackage{dvipsnames,svgnames,x11names}{xcolor}
\documentclass[
  a4paper,
]{article}
\usepackage{xcolor}
\usepackage[margin=25mm]{geometry}
\usepackage{amsmath,amssymb}
\setcounter{secnumdepth}{-\maxdimen} % remove section numbering
\usepackage{iftex}
\ifPDFTeX
  \usepackage[T1]{fontenc}
  \usepackage[utf8]{inputenc}
  \usepackage{textcomp} % provide euro and other symbols
\else % if luatex or xetex
  \usepackage{unicode-math} % this also loads fontspec
  \defaultfontfeatures{Scale=MatchLowercase}
  \defaultfontfeatures[\rmfamily]{Ligatures=TeX,Scale=1}
\fi
\usepackage{lmodern}
\ifPDFTeX\else
  % xetex/luatex font selection
  \setmainfont[]{DejaVu Serif}
  \setsansfont[]{DejaVu Sans}
  \setmonofont[]{DejaVu Sans Mono}
\fi
% Use upquote if available, for straight quotes in verbatim environments
\IfFileExists{upquote.sty}{\usepackage{upquote}}{}
\IfFileExists{microtype.sty}{% use microtype if available
  \usepackage[]{microtype}
  \UseMicrotypeSet[protrusion]{basicmath} % disable protrusion for tt fonts
}{}
\makeatletter
\@ifundefined{KOMAClassName}{% if non-KOMA class
  \IfFileExists{parskip.sty}{%
    \usepackage{parskip}
  }{% else
    \setlength{\parindent}{0pt}
    \setlength{\parskip}{6pt plus 2pt minus 1pt}}
}{% if KOMA class
  \KOMAoptions{parskip=half}}
\makeatother
\usepackage{color}
\usepackage{fancyvrb}
\newcommand{\VerbBar}{|}
\newcommand{\VERB}{\Verb[commandchars=\\\{\}]}
\DefineVerbatimEnvironment{Highlighting}{Verbatim}{commandchars=\\\{\}}
% Add ',fontsize=\small' for more characters per line
\newenvironment{Shaded}{}{}
\newcommand{\AlertTok}[1]{\textcolor[rgb]{1.00,0.00,0.00}{\textbf{#1}}}
\newcommand{\AnnotationTok}[1]{\textcolor[rgb]{0.38,0.63,0.69}{\textbf{\textit{#1}}}}
\newcommand{\AttributeTok}[1]{\textcolor[rgb]{0.49,0.56,0.16}{#1}}
\newcommand{\BaseNTok}[1]{\textcolor[rgb]{0.25,0.63,0.44}{#1}}
\newcommand{\BuiltInTok}[1]{\textcolor[rgb]{0.00,0.50,0.00}{#1}}
\newcommand{\CharTok}[1]{\textcolor[rgb]{0.25,0.44,0.63}{#1}}
\newcommand{\CommentTok}[1]{\textcolor[rgb]{0.38,0.63,0.69}{\textit{#1}}}
\newcommand{\CommentVarTok}[1]{\textcolor[rgb]{0.38,0.63,0.69}{\textbf{\textit{#1}}}}
\newcommand{\ConstantTok}[1]{\textcolor[rgb]{0.53,0.00,0.00}{#1}}
\newcommand{\ControlFlowTok}[1]{\textcolor[rgb]{0.00,0.44,0.13}{\textbf{#1}}}
\newcommand{\DataTypeTok}[1]{\textcolor[rgb]{0.56,0.13,0.00}{#1}}
\newcommand{\DecValTok}[1]{\textcolor[rgb]{0.25,0.63,0.44}{#1}}
\newcommand{\DocumentationTok}[1]{\textcolor[rgb]{0.73,0.13,0.13}{\textit{#1}}}
\newcommand{\ErrorTok}[1]{\textcolor[rgb]{1.00,0.00,0.00}{\textbf{#1}}}
\newcommand{\ExtensionTok}[1]{#1}
\newcommand{\FloatTok}[1]{\textcolor[rgb]{0.25,0.63,0.44}{#1}}
\newcommand{\FunctionTok}[1]{\textcolor[rgb]{0.02,0.16,0.49}{#1}}
\newcommand{\ImportTok}[1]{\textcolor[rgb]{0.00,0.50,0.00}{\textbf{#1}}}
\newcommand{\InformationTok}[1]{\textcolor[rgb]{0.38,0.63,0.69}{\textbf{\textit{#1}}}}
\newcommand{\KeywordTok}[1]{\textcolor[rgb]{0.00,0.44,0.13}{\textbf{#1}}}
\newcommand{\NormalTok}[1]{#1}
\newcommand{\OperatorTok}[1]{\textcolor[rgb]{0.40,0.40,0.40}{#1}}
\newcommand{\OtherTok}[1]{\textcolor[rgb]{0.00,0.44,0.13}{#1}}
\newcommand{\PreprocessorTok}[1]{\textcolor[rgb]{0.74,0.48,0.00}{#1}}
\newcommand{\RegionMarkerTok}[1]{#1}
\newcommand{\SpecialCharTok}[1]{\textcolor[rgb]{0.25,0.44,0.63}{#1}}
\newcommand{\SpecialStringTok}[1]{\textcolor[rgb]{0.73,0.40,0.53}{#1}}
\newcommand{\StringTok}[1]{\textcolor[rgb]{0.25,0.44,0.63}{#1}}
\newcommand{\VariableTok}[1]{\textcolor[rgb]{0.10,0.09,0.49}{#1}}
\newcommand{\VerbatimStringTok}[1]{\textcolor[rgb]{0.25,0.44,0.63}{#1}}
\newcommand{\WarningTok}[1]{\textcolor[rgb]{0.38,0.63,0.69}{\textbf{\textit{#1}}}}
\usepackage{longtable,booktabs,array}
\newcounter{none} % for unnumbered tables
\usepackage{calc} % for calculating minipage widths
% Correct order of tables after \paragraph or \subparagraph
\usepackage{etoolbox}
\makeatletter
\patchcmd\longtable{\par}{\if@noskipsec\mbox{}\fi\par}{}{}
\makeatother
% Allow footnotes in longtable head/foot
\IfFileExists{footnotehyper.sty}{\usepackage{footnotehyper}}{\usepackage{footnote}}
\makesavenoteenv{longtable}
\setlength{\emergencystretch}{3em} % prevent overfull lines
\providecommand{\tightlist}{%
  \setlength{\itemsep}{0pt}\setlength{\parskip}{0pt}}
\usepackage{luatexja}
\usepackage{xurl}
\usepackage{bookmark}
\IfFileExists{xurl.sty}{\usepackage{xurl}}{} % add URL line breaks if available
\urlstyle{same}
\hypersetup{
  pdftitle={The Fourteenth Thought Experiment: Why Can Different Symmetries Never Be Fully Reconciled?},
  pdfauthor={Noriaki Kihara (WF System Co., Ltd.) ORCID 0009-0004-6753-4020},
  colorlinks=true,
  linkcolor={Maroon},
  filecolor={Maroon},
  citecolor={Blue},
  urlcolor={Blue},
  pdfcreator={LaTeX via pandoc}}

\title{The Fourteenth Thought Experiment: Why Can Different Symmetries
Never Be Fully Reconciled?}
\usepackage{etoolbox}
\makeatletter
\providecommand{\subtitle}[1]{% add subtitle to \maketitle
  \apptocmd{\@title}{\par {\large #1 \par}}{}{}
}
\makeatother
\subtitle{― Non-commensurability, Pull-in and Metastability of
Integer-Selecting Symmetries, and Readout along the Direction of
Evolution ―}
\author{Noriaki Kihara (WF System Co., Ltd.) ORCID 0009-0004-6753-4020}
\date{October 12, 2026 Version DOI 10.5281/zenodo.23315558}

\ltjsetparameter{jacharrange={-2,-3}}
\begin{document}
\maketitle

\textbf{Author:} Noriaki Kihara\\
\textbf{ORCID:} 0009-0004-6753-4020\\
\textbf{Version DOI:} 10.5281/zenodo.23315558\\
\textbf{Concept DOI:} 10.5281/zenodo.23315557\\
\textbf{Version:} v1.0 (a discovery record from dialogue; the research
hypothesis fixed in writing)\\
\textbf{Date:} 2026-10-12\\
\textbf{Related papers:} Twelfth Thought Experiment {[}K6{]}, Thirteenth
Thought Experiment {[}K7{]}, finite-order resonance {[}K4{]}\\
\textbf{Classification:} thought experiment / exploratory theoretical
conception / verification tasks for numerical correspondences\\
\textbf{Recording method:} Starting from dialogues between the author
and AI (first draft with ChatGPT, revision with Claude), organized while
preserving the order of the author's corrections and questions.
Recalculation of numbers and checking of related literature were done as
support.\\
\textbf{Keywords:} discrete symmetry, competing symmetries, integer
selection, number-theoretic selection, non-commensurability, pull-in,
incomplete locking, phase slip, metastability, direction of evolution,
emergence of time, origin of uncertainty, projective readout,
frustration, E₈, zero square sum, central projection, anonymity

\begin{center}\rule{0.5\linewidth}{0.5pt}\end{center}

\subsection{Abstract}\label{abstract}

This paper began with a dialogue about the quantum-mechanical
explanation of ``why the electron does not fall into the nucleus''.
Planck's constant and the uncertainty principle explain the stability of
the hydrogen atom, but they do not ask why quantum structure itself
exists. The author focused on why the Standard Model requires several
heterogeneous symmetries, on the difference between the integer
structure of charge and the non-integer hierarchy of masses, and on the
discrete selection conditions for space, time and causal order.

The central hypothesis of this paper is as follows. \textbf{Several
symmetries that select integers, in space, in the order of evolution, in
charge and elsewhere, are required simultaneously on one relational
structure; they cannot be satisfied completely at the same time, yet
they pull each other in. This incomplete locking produces metastability
and sudden change, and the observed constants appear as effective
quantities read out from that history along a particular direction.} The
author's outlook that, without symmetries requiring integers for order
and causality, neither time nor uncertainty would arise is also part of
this hypothesis.

\textbf{The proposal of this paper is to specify what should be searched
for under this hypothesis.} That is: \textbf{search for pairs that
appear theoretically unrelated, with neither an exact map nor a line of
derivation between them, and yet lie close to each other in number or
structure. Such a pair pulls in when the two are imposed together, but
never locks completely.} There is no need to explain why they are close.
Being near but not identical is itself the clue for finding candidates
for incomplete locking, and a pair for which an exact agreement (a lock)
is found is removed from the candidates.

What led the author to notice this structure was the fine-structure
constant. The four-dimensional cell count \(137=1+8+128\) and the
finite-order root of two-wave exchange \(R_{124,23}=\cos^2(23\pi/124)\)
each give, independently, a value near \(\alpha^{-1}\). This paper
treats this in Chapter 7 as \textbf{one case} of a pair that is near but
not identical. The numerical closeness is treated not as a proof of
agreement but as a clue for the search, and it is recorded that the
four-dimensional 137 and 23/124 also share integers with E₈ while
differing in form. Further, above these integer-selecting symmetries the
author places the zero square sum \(\sum x^2=0\), which the author
positions as the ultimate symmetry of the author's theoretical system.
Chapter 11 states that, like \(U^n=I\), it holds not as an assumption
but as a consequence of combining other symmetries; that it is a
definition expressing that a conserved quantity can exist, rather than
fixing the value of a conserved quantity to 0 (the value \(R\) moves to
an axis \(iR\)); that central projection separates observable from
unobservable directions; the anonymity of axes; and the sharing of the
origin under constant curvature. This paper is not one that completes
the derivation of a specific constant; it is a discovery record that
fixes a general principle and verification tasks.

\begin{center}\rule{0.5\linewidth}{0.5pt}\end{center}

\subsection{0. Research Stance Fixed First in This
Paper}\label{research-stance-fixed-first-in-this-paper}

\begin{enumerate}
\def\labelenumi{\arabic{enumi}.}
\tightlist
\item
  \textbf{Keep the order of questions.} Symmetry comes first; conserved
  quantities are consequences that can be read by projecting onto
  particular directions. Do not go straight from continuous symmetries
  to conserved quantities.
\item
  \textbf{Do not give \(U^n=I\) and \(\sum x^2=0\) in advance as
  explanatory principles.} Finite recurrence and the zero square sum are
  properties observed after examining interactions and integer
  conditions, or consequences of combining several symmetries (§11.1).
\item
  \textbf{Searching only for a ``common fixed point'' of several
  symmetries is not enough.} The idea of this paper concerns the case in
  which different conditions cannot be fully compatible and nevertheless
  pull each other in.
\item
  \textbf{Do not posit physical time a priori.} The index of evolution
  is a computational order; the readout condition that identifies it
  with physical time is a separate problem.
\item
  \textbf{Do not fit observed values afterwards to the integers treated
  as cases.} The reasons for selecting the integers and the maps to the
  observed constants are to be specified and verified independently of
  the respective observed values.
\item
  \textbf{Do not add a third symmetry by assumption.} Consider a missing
  constraint only when two conditions are not enough.
\end{enumerate}

\subsection{1. How the Idea Arose: The Order of
Questions}\label{how-the-idea-arose-the-order-of-questions}

The idea of this paper took shape along the order of the author's
corrections in the dialogue (the original statements are in Chapter 9).

{\def\LTcaptype{none} % do not increment counter
\begin{longtable}[]{@{}
  >{\raggedright\arraybackslash}p{(\linewidth - 4\tabcolsep) * \real{0.3333}}
  >{\raggedright\arraybackslash}p{(\linewidth - 4\tabcolsep) * \real{0.3333}}
  >{\raggedright\arraybackslash}p{(\linewidth - 4\tabcolsep) * \real{0.3333}}@{}}
\toprule\noalign{}
\begin{minipage}[b]{\linewidth}\raggedright
Stage
\end{minipage} & \begin{minipage}[b]{\linewidth}\raggedright
Author's point (statement number in Chapter 9)
\end{minipage} & \begin{minipage}[b]{\linewidth}\raggedright
Question fixed at that stage
\end{minipage} \\
\midrule\noalign{}
\endhead
\bottomrule\noalign{}
\endlastfoot
2026-10-10 & In the Thirteenth Thought Experiment, dissipation,
reversible exchange and particle lifetime were examined, and it was
stated that ``\(U^n=I\) is a result, not a cause'' & The mechanism by
which two or more symmetries select long-lived states remains
unsolved \\
Dialogue 1 & Planck's constant is not the end point of explanation. The
structure that determines the symmetries of spacetime is fundamental
(statement 1) & Ask about the origin of quantum structure, going back
before \(\hbar\) \\
Dialogue 2 & Combinations of several symmetries require discreteness and
periodicity of space and time, and conserved quantities are read by
projection onto particular directions (statement 2) & Start not from
conserved quantities but from symmetries \\
Dialogue 3 & Discretization requires symmetries that give integers
(magic numbers). Order and causality also require symmetries demanding
integers; without them neither time nor uncertainty arises (statements 3
and 4) & Ask about integer selection not only on the side of space but
also on the side of order and causality \\
Dialogue 4 & From here on it is not the world of physical symmetries but
of number theory (statement 5) & Treat the problem of symmetry as a
problem of selecting sets of integers \\
Dialogue 5 & Several symmetries cannot hold simultaneously; they pull in
but never lock completely, and this produces sudden change and
metastability (statement 6) & Search not for common solutions but for
combinations of non-commensurability and pull-in \\
Dialogue 6 & What led to the noticing was the existence of two 137s near
the fine-structure constant (statements 7--9) & Examine the
four-dimensional 137 and the 23/124 of the direction of evolution as a
case of the general principle \\
Revision dialogue & α is one case. Precedents for integer selection and
E₈. The zero square sum is a consequence, not an assumption, and a
definition that a conserved quantity can exist. Being near but not
identical is the content of the hypothesis (statements 12--20) & Search
for pairs that appear theoretically unrelated yet lie close, pull in,
but never lock completely \\
\end{longtable}
}

Therefore, the new point of this paper is not the discovery of specific
numbers but \textbf{the reinterpretation of the fact that different
integer-selecting structures cannot be fully reconciled as a candidate
condition that, on the contrary, produces emergent oscillation,
stability and readout}, and the proposal of this as a search strategy:
\textbf{search for pairs that appear theoretically unrelated yet lie
close, and that pull in but never lock completely}.

\subsection{2. The First Question: Is Planck's Constant the End Point or
the Starting Point of
Explanation?}\label{the-first-question-is-plancks-constant-the-end-point-or-the-starting-point-of-explanation}

The energy of a hydrogen-like atom has, non-relativistically, a finite
minimum from the competition

\[
E(r)\sim\frac{\hbar^2}{2\mu r^2}-\frac{e^2}{4\pi\varepsilon_0r}
\]

This is a valid explanation in standard quantum mechanics. This
calculation, however, takes \(\hbar\) and the quantum structure of
position and momentum as input. The author's question is \textbf{why
such quantum structures become necessary}. This paper therefore does not
deny the results of standard theory, but asks about the origin of the
structure, going further back.

Instead of giving \(\hbar\) and coupling constants first as observed
quantities, it asks whether they can be reconstructed from relations,
symmetries, integer selection, stability and observation maps. Merely
rewriting the equations with \(\hbar\) removed, however, is not a
derivation of \(\hbar\).

\subsection{3. What Should Be Searched For: The Simultaneous Requirement
of Heterogeneous Integer-Selecting
Symmetries}\label{what-should-be-searched-for-the-simultaneous-requirement-of-heterogeneous-integer-selecting-symmetries}

\subsubsection{3.1 Heterogeneous
Symmetries}\label{heterogeneous-symmetries}

The symmetries of the Standard Model have a spacetime side (Lorentz
transformations) and an internal side
(\(SU(3)_C\times SU(2)_L\times U(1)_Y\)), and in addition there are
unexplained integer and algebraic structures such as the number of
generations and the mass hierarchy. The author's claim is not to list
what each symmetry conserves, but to ask \textbf{why several conditions
of such different character are required at the same time}. In the
author's theoretical system, the symmetry of the zero square sum
\(\sum x^2=0\), which has neither a scale nor names for the axes, is
placed above these (Chapter 11).

\subsubsection{3.2 Symmetries That Give Integers (Magic
Numbers)}\label{symmetries-that-give-integers-magic-numbers}

Continuous symmetries alone do not produce discretization. For
discretization to arise, symmetries are needed that give specific
integers to the allowed numbers of configurations, directions and
recurrences. In this paper these are called ``integer-selecting
symmetries''.

\subsubsection{3.3 Discretization of Order and Causality: Time and
Uncertainty}\label{discretization-of-order-and-causality-time-and-uncertainty}

Integer selection on the side of space alone is not enough. Integers
giving a minimal distinguishable structure are also required for the
order of evolution and for causal relations. In the author's outlook,
without such symmetries neither time nor uncertainty would arise. There
are three footholds from known structures for this outlook.

\textbf{(a) The scale of time can be made as an integer count of beats.}
To make a scale of time without positing an external clock, it suffices
that two mutually coherent oscillations with different frequencies
coexist. Each time their relative phase shifts by \(2\pi\) (each time
one of them gains a full turn on the other), the intensity of the sum
returns to the same brightness. This single shift is an \textbf{integer
event} without fractions, and the scale is counted as the number of such
events. Clock comparison by heterodyne detection and optical frequency
combs measure time by this principle. Duration is read not as a length
of delay but as a number of integer slips. The phase slip of H4 in
Chapter 4 is an integer event of the same kind.

\textbf{(b) The cost of having a scale becomes the frame of
uncertainty.} To mark a scale with spacing \(\delta t\), a coherent pair
with frequency difference \(\Delta\nu\gtrsim1/\delta t\) is needed.
Hence

\[
\delta t\,\Delta\nu\gtrsim1
\]

appears not as an axiom imposed from outside but as the cost of making
the scale (Gabor's time-frequency cell {[}12{]}). What comes out here is
the \textbf{form} of the inequality; no value linking its unit to
Planck's constant \(h\) comes out. The derivation of the value remains
with the question of §2.

\textbf{(c) The phases that can be marked exactly with integers are
limited.} If a phase is represented by a pair of integers \((a,b)\) and
its rotation is carried out exactly with integer operations alone, the
allowed steps are limited, for the same reason as in the
crystallographic restriction theorem, to the two families of 90°
(\(Z_4\)) and 60° (\(Z_6\)) (§10.1.1). Integer selection on the side of
evolution comes from the same theorem as the crystallographic
restriction of space.

Using (a)--(c) as footholds, this paper fixes \textbf{the question that
seeks the origin of time and uncertainty in the integer selection of
order and causality}. The starting point for thinking about uncertainty
without taking \(\hbar\) as input lies here.

\subsubsection{3.4 Transition to Number
Theory}\label{transition-to-number-theory}

At this point the problem becomes not the properties of individual
physical symmetries but \textbf{the problem of characterizing the sets
of integers simultaneously allowed by several symmetries}. As
mathematical precedents there are the crystallographic restriction
theorem, Hurwitz's dimension restriction and E₈ (Chapter 10).

\subsubsection{3.5 Candidates for Integer
Selection}\label{candidates-for-integer-selection}

The integer structures to be selected are not limited to one. Candidates
to which the hypothesis of this paper may apply are listed.

{\def\LTcaptype{none} % do not increment counter
\begin{longtable}[]{@{}
  >{\raggedright\arraybackslash}p{(\linewidth - 4\tabcolsep) * \real{0.3333}}
  >{\raggedright\arraybackslash}p{(\linewidth - 4\tabcolsep) * \real{0.3333}}
  >{\raggedright\arraybackslash}p{(\linewidth - 4\tabcolsep) * \real{0.3333}}@{}}
\toprule\noalign{}
\begin{minipage}[b]{\linewidth}\raggedright
Candidate
\end{minipage} & \begin{minipage}[b]{\linewidth}\raggedright
Example of selected integers
\end{minipage} & \begin{minipage}[b]{\linewidth}\raggedright
Treatment in this paper
\end{minipage} \\
\midrule\noalign{}
\endhead
\bottomrule\noalign{}
\endlastfoot
Number of directions read as space & Number of dimensions, and why its
increase is not allowed without limit & Task: identify the two
constraints that select this integer, in the form of §4 \\
Minimal number of structures distinguishing the order and causality of
evolution & Integer slips of beats (§3.3(a)) & Fixed as a question in
§3.3 \\
Phases that can be marked exactly with integers & Orders 1, 2, 3, 4, 6
(\(Z_4\), \(Z_6\)) & §3.3(c), §10.1.1 \\
The difference between charge quantum numbers and non-integer masses and
couplings & \(q=nq_0,\ q_0=e/3\) & Task: identify the two constraints
that select this integer, in the form of §4 \\
Number of generations and mass hierarchy & 3 generations & Task:
identify the two constraints that select this integer, in the form of
§4 \\
Finite recurrence of phase, resonance, metastability, discrete
transitions & Finite-order root \(R_{n,m}\) & Part of the case in
Chapter 7 \\
Lattices satisfying evenness and self-duality at the same time &
Dimension 8 (E₈), 240 roots & §10.3 \\
Integer count of four-dimensional geometry and finite-order root of the
direction of evolution & 137, 23/124 & \textbf{Case in Chapter 7
(fine-structure constant)} \\
\end{longtable}
}

\subsection{4. The Central Hypothesis of This Paper: Symmetries That
Cannot Be Fully Reconciled and Still Pull
In}\label{the-central-hypothesis-of-this-paper-symmetries-that-cannot-be-fully-reconciled-and-still-pull-in}

Call the two families of constraints, provisionally,
\(\mathcal C_{\rm A}\) (the side of space) and \(\mathcal C_{\rm B}\)
(the side of evolution).

If the two conditions simply have common solutions, a completely fixed
state may be explained, but metastable oscillation and sudden switching
are not necessarily explained. What is sought here is a structure that
has the following properties at the same time.

\textbf{H1 (selection):} The constraint of each symmetry imposes integer
restrictions on the allowed numbers of configurations, numbers of
directions and order relations.\\
\textbf{H2 (non-commensurability):} No physical state satisfies the two
constraints completely at the same time, or such a state cannot persist
under evolution.\\
\textbf{H3 (pull-in):} The interaction tends to reduce the deviation
from each constraint and creates locally long residence times.\\
\textbf{H4 (incomplete locking):} Complete fixing is not reached; phase
slips or state transitions occur occasionally.\\
\textbf{H5 (readout):} An integral readout of the evolution gives an
effective constant that depends on the measurement conditions.

\textbf{The search proposal (the core of this paper's claim).}
Structures satisfying H1--H5 are not derived from theory but searched
for by the following four conditions.

\begin{enumerate}
\def\labelenumi{\arabic{enumi}.}
\tightlist
\item
  \textbf{Appear theoretically unrelated:} neither an exact map nor a
  line of derivation between the two structures is known.
\item
  \textbf{Nevertheless lie close:} they share integers or values, or
  their values are close. But their forms or values do not agree.
\item
  \textbf{Pull in:} when the two are imposed together, the state moves
  toward each side and stays locally long (H3).
\item
  \textbf{Never lock completely:} agreement is not reached; phase slips,
  metastability and sudden change remain (H2, H4).
\end{enumerate}

The reason for the closeness is not explained. Closeness is an input to
the search, not a result to be explained. A pair for which an exact
correspondence is found is locked and is removed from the candidates; a
pair that is completely unrelated does not pull in and is not a
candidate. Known examples are the edges of the tongues of the circle map
(end of this chapter), the quasicrystals that arise from incompatible
symmetries (§10.1), and the case of Chapter 7 (the four-dimensional 137,
23/124, E₈ and the values near the fine-structure constant).

In an abstract model for verification, one may define two non-negative
violation measures \(D_{\rm A}(S),D_{\rm B}(S)\) for a state \(S\) and
consider the case

\[
\operatorname*{argmin}D_{\rm A}
\;\cap\;
\operatorname*{argmin}D_{\rm B}
=\varnothing
\tag{1}
\]

Constructing \(D_{\rm A},D_{\rm B}\) satisfying Eq. (1) from a concrete
update rule is Task 1. The evaluation functions are readouts for
analyzing the state, not internal variables to be added to the update
rule without permission.

\textbf{From a mere mismatch of integers, such as the integers selected
by the two constraints being coprime, complete non-commensurability
cannot be concluded.} If both are regarded simply as periods, a least
common multiple exists. To prove the competition, one must show the
operational conditions that both sides impose on the same state.

\textbf{A known example: the circle map.} The simplest known example
satisfying H1--H4 at the same time is the circle map

\[
\theta_{k+1}=\theta_k+\Omega-\frac{K}{2\pi}\sin2\pi\theta_k\pmod 1
\]

This is an update rule containing no external time, of the form of Eq.
(2) in Chapter 5. For \(0<K<1\), intervals of \(\Omega\) open (Arnold
tongues) in which the mean rotation number \(\rho\) is pulled in to a
rational number \(p/q\), and a pair of integers \((p,q)\) is selected
(H1). The driving rotation \(\Omega\) and the rational rotation favored
by the nonlinear term generally do not agree (H2). Inside a tongue the
phase is pulled in to \(p/q\) (H3). Just outside the edge of a tongue,
long nearly locked residences and short slips in which the phase shifts
by one full turn alternate (H4). At \(K=1\) the pull-in intervals occupy
the full measure of \(\Omega\), and \(\rho(\Omega)\) becomes a devil's
staircase {[}11{]}. The hypothesis of this paper can be read as
replacing the two periods given from outside in this example (\(\Omega\)
and the period of the nonlinear term) by two integer selections arising
from inside the relational structure.

\subsection{5. Reading along the Direction of Evolution without Assuming
External
Time}\label{reading-along-the-direction-of-evolution-without-assuming-external-time}

Write the basic update as

\[
S_{k+1}=F(S_k)
\tag{2}
\]

Here \(k\) is a computational update number, not a physical time \(t\).
The requirement is kept that the update rule \(F\) needs no hidden clock
or event label other than the state.

If the readout of the observed interaction amplitude is \(A_q(S_k)\),
one can formally set, as a candidate for the quantity obtained from the
history along the direction of evolution,

\[
\overline{A}_q
=\lim_{N\to\infty}
\frac{\sum_{k=0}^{N-1}w_q(k)\,A_q(S_k)}
{\sum_{k=0}^{N-1}w_q(k)}
\tag{3}
\]

Here \(q\) is an observational scale defined afterwards, and \(w_q\) is
the weight required by the measurement response. On that basis, the
effective constant that is read out is taken as

\[
c_{\mathrm{read}}(q)=\mathcal R_q(\overline A_q)
\tag{4}
\]

The fine-structure constant \(\alpha_{\rm read}(q)\) is one example of
it (Chapter 7).

\textbf{This is not a claim that simply averaging instantaneous values
gives the observed value.} The readout after integrating complex
amplitudes including phase generally differs from the average of the
constant itself. Moreover, if \(w_q\) or \(\mathcal R_q\) could be
adjusted freely afterwards, any number could be obtained; they too must
be defined independently from the interaction and the observational
structure.

Only if metastable oscillation, sudden switching and finite lifetimes
appear as \emph{results} of Eq. (2), and the response can be read as Eq.
(4), does the central hypothesis of this paper acquire dynamical
content.

\subsection{6. Verification Tasks and Falsification Conditions for the
General
Principle}\label{verification-tasks-and-falsification-conditions-for-the-general-principle}

The priority is not the search for solutions of \(U^n=I\) but
\textbf{the discovery of integer-selecting constraints and their
competition}.

\textbf{Task 1: Proof of non-commensurability.} State explicitly in
which operation the two constraints collide on the same state. If common
solutions exist in the ordinary way, the hypothesis of ``impossibility
of complete locking'' is rejected.

\textbf{Task 2: Metastable dynamics.} Reproduce distributions of
residence times, frequencies of phase slips, finite lifetimes and sudden
changes without hidden clocks, additional state variables or post-hoc
selection rules. A third constraint is not added until its necessity is
shown.

\textbf{Task 3: Observation map.} Define \(c_{\rm read}(q)\) from
interaction amplitudes and predict values at different scales with a
common parameter setting. Fix the definition before any numerical
fitting.

\textbf{Task 4: Integer selection of order and causality.} Show whether
the minimal structure distinguishing the order and causality of
evolution requires integers, and how this appears in the readout of time
and uncertainty (§3.3). As a concrete candidate, there is the
combination that assigns to \(\mathcal C_{\rm A}\) the rotation
restriction of spatial lattices (§10.1) and to \(\mathcal C_{\rm B}\)
the restriction of phases that can be marked with integers (§10.1.1).
The rotation orders allowed by four-dimensional spatial lattices,
\(\{1,2,3,4,5,6,8,10,12\}\), and the orders allowed for phases
represented by a pair of integers, \(\{1,2,3,4,6\}\), do not agree.
Examine whether this mismatch becomes a collision of operations on the
same state (H2).

If the necessary conditions fail in any of these, the strong claim of
this paper is rejected or revised.

\subsection{7. Case: The Four-Dimensional 137 and the 23/124 of the
Direction of Evolution (Fine-Structure
Constant)}\label{case-the-four-dimensional-137-and-the-23124-of-the-direction-of-evolution-fine-structure-constant}

This chapter treats the case that led the author to notice this
structure. The verification of the general hypothesis is in Chapter 6;
the numbers in this chapter remain one example of it.

\subsubsection{7.1 Research History of the
Case}\label{research-history-of-the-case}

{\def\LTcaptype{none} % do not increment counter
\begin{longtable}[]{@{}
  >{\raggedright\arraybackslash}p{(\linewidth - 4\tabcolsep) * \real{0.3333}}
  >{\raggedright\arraybackslash}p{(\linewidth - 4\tabcolsep) * \real{0.3333}}
  >{\raggedright\arraybackslash}p{(\linewidth - 4\tabcolsep) * \real{0.3333}}@{}}
\toprule\noalign{}
\begin{minipage}[b]{\linewidth}\raggedright
Date
\end{minipage} & \begin{minipage}[b]{\linewidth}\raggedright
Research record / result
\end{minipage} & \begin{minipage}[b]{\linewidth}\raggedright
Problem remaining at that time
\end{minipage} \\
\midrule\noalign{}
\endhead
\bottomrule\noalign{}
\endlastfoot
2026-04-21 & Conceived the integer representation \(137=8^2+8^2+3^2\) &
The mechanism by which it becomes an actual physical constant is
unknown \\
2026-05-14 & Examined the energy dependence of \(\alpha\) and the
hypothesis of 137 as a ground structure & The connection between crystal
geometry and the running in the Standard Model is incomplete \\
2026-05-30 & Organized the number 137 of unit cells inside a
four-dimensional ball of radius 3 (the three-layer representation at
that time was 9+64+64) & Why four dimensions and radius 3 are selected
is not derived \\
2026-07-17 & Wrote up the four-dimensional cell count \(137=1+8+128\),
the geometric self-consistent value and the readout of the 128 outer
shell as a paper & The correspondence to electromagnetic coupling and an
independent selection principle are unproven \\
2026-07 to 09 & Analyzed and numerically verified the finite-order root
23/124 of the two-wave exchange operator & The periodicity obtained by
giving \(R\) must be distinguished from a mechanism by which \(R\) is
selected spontaneously \\
\end{longtable}
}

\subsubsection{7.2 Integer Structure on the Side of Space: The
Four-Dimensional 137
Cells}\label{integer-structure-on-the-side-of-space-the-four-dimensional-137-cells}

The author's earlier work used the condition that a four-dimensional
cube of side 1 with integer center
\(c=(c_1,c_2,c_3,c_4)\in\mathbb Z^4\),

\[
C_c=c+[-1/2,1/2]^4
\]

is completely contained in the four-dimensional ball of radius 3
centered at the origin:

\[
C_c\subseteq B_4(3)
\quad\Longleftrightarrow\quad
\sum_{j=1}^4\left(|c_j|+\frac12\right)^2\le9.
\tag{5}
\]

The number of cells satisfying this condition is, by direct enumeration,
exactly

\[
\boxed{N_4(3)=137=1+8+128}
\tag{6}
\]

This is a division into one central cell, eight axial neighbouring cells
and 128 other cells. The attached re-verification script performs this
count independently.

Here, \textbf{that a cell count of 137 is obtained} and \textbf{that
nature spontaneously selects this sector of four dimensions and radius
3} are entirely different. The latter is unproven and is included in the
task of integer selection asked in §3.

From the geometric self-consistency relation of the earlier work,

\[
\alpha_g^{-1}=137+\frac{\pi^2}{2}\alpha_g
\tag{7}
\]

setting \(x=\alpha_g^{-1}\) gives

\[
x_g=\frac{137+\sqrt{137^2+2\pi^2}}{2}
=137.036010988388\ldots
\tag{8}
\]

Eq. (7), however, is a relation proposed in the earlier work, not a
physical law necessarily derived from the lattice count alone. Below,
\(\alpha_g\) is distinguished from the observed electromagnetic coupling
\(\alpha_{\mathrm{EM}}\).

\subsubsection{7.3 Integer Structure on the Side of Evolution: The
Finite-Order Root
23/124}\label{integer-structure-on-the-side-of-evolution-the-finite-order-root-23124}

Let the finite-order candidates for the reflectance in the earlier
two-channel exchange experiment be

\[
R_{n,m}=\cos^2\left(\frac{\pi m}{n}\right),
\qquad \gcd(m,n)=1
\tag{9}
\]

For \(n=124,m=23\),

\[
R_{124,23}=0.6971779275566594\ldots
\tag{10}
\]

and the eigenvalues of the symmetric and antisymmetric modes of the
exchange operator show finite recurrence. This is an analytic
consequence of giving the existing operator and Eq. (9).

The correspondence candidate adopted in the earlier work,

\[
x_{23/124}=\frac{4\pi}{(1-R_{124,23})^2}
=\frac{4\pi}{\sin^4(23\pi/124)}
=137.036042914599\ldots
\tag{11}
\]

lies near the observed value.

\textbf{However, neither that Eq. (11) physically represents
\(\alpha^{-1}\) nor that 23/124 is spontaneously selected from the
infinitely many finite-order roots has been derived.} As emphasized in
the Thirteenth Thought Experiment, one must not place finite recurrence
\(U^n=I\) first as a cause and regard stability as thereby explained.

\subsubsection{7.4 Do Not Confuse Independent Numbers and
Residuals}\label{do-not-confuse-independent-numbers-and-residuals}

{\def\LTcaptype{none} % do not increment counter
\begin{longtable}[]{@{}
  >{\raggedright\arraybackslash}p{(\linewidth - 4\tabcolsep) * \real{0.3000}}
  >{\raggedleft\arraybackslash}p{(\linewidth - 4\tabcolsep) * \real{0.4000}}
  >{\raggedright\arraybackslash}p{(\linewidth - 4\tabcolsep) * \real{0.3000}}@{}}
\toprule\noalign{}
\begin{minipage}[b]{\linewidth}\raggedright
Object
\end{minipage} & \begin{minipage}[b]{\linewidth}\raggedleft
Value proposed to correspond to \(\alpha^{-1}\)
\end{minipage} & \begin{minipage}[b]{\linewidth}\raggedright
Mathematical / observational status
\end{minipage} \\
\midrule\noalign{}
\endhead
\bottomrule\noalign{}
\endlastfoot
Four-dimensional cell count & \(137\) & Exact integer count under the
specified geometric condition \\
Geometric self-consistency \(x_g\) & \(137.036010988389\ldots\) &
Mathematical solution of the proposed self-consistency equation \\
Finite-order-root correspondence \(x_{23/124}\) &
\(137.036042914599\ldots\) & Value of the proposed correspondence
formula from \(R\) \\
CODATA 2022 \(\alpha_{\mathrm{EM}}^{-1}(0)\) & \(137.035999177(21)\) &
Recommended observed value in the low-momentum limit \\
\end{longtable}
}

The computed differences are

\[
\begin{aligned}
x_{23/124}-x_g&=0.000031926210\ldots,\\
x_g-\alpha_{\mathrm{EM}}^{-1}(0)&=0.000011811389\ldots,\\
x_{23/124}-\alpha_{\mathrm{EM}}^{-1}(0)&=0.000043737599\ldots.
\end{aligned}
\tag{12}
\]

These differences greatly exceed the measurement uncertainty.
\textbf{Numerical closeness does not prove a common physical origin.} In
the hypothesis of this paper, the very fact that the three values are
near but not identical may be a manifestation of several structures
pulling each other in without locking completely (§7.6). Searching for a
correction mechanism that fills the differences is not a task of this
paper.

\subsubsection{7.5 Application to the General
Hypothesis}\label{application-to-the-general-hypothesis}

In this case, the families of constraints of Chapter 4 are examined as
candidates assigning \(\mathcal C_{\rm A}\) to the 137 of
four-dimensional geometry and \(\mathcal C_{\rm B}\) to the 23/124 of
the direction of evolution. Complete non-commensurability cannot be
concluded from the fact that 137 and 124 are coprime: if both are
regarded simply as periods, the least common multiple 16988 exists. Only
when the operational conditions that both sides impose on the same state
are shown does this case satisfy H2 of Chapter 4.

If this assignment holds and the value obtained by reading the history
of incomplete locking with Eq. (4) gives \(\alpha_{\rm read}(q)\), the
observed fine-structure constant appears as an effective quantity
composed of two or more discrete constraints.

\subsubsection{7.6 Integer Correspondences with
E₈}\label{integer-correspondences-with-eux2088}

E₈, treated in Chapter 10 (§10.3), shares the following integers with
the two integer structures of this chapter. They are recorded not in the
order in which the correspondences were made but as integer relations
that can be checked (all items can be recalculated with the attached
\texttt{verify\_\allowbreak{}e8.py}).

\textbf{Side of space: 137 and 248.}

\[
137=1+8+128,\qquad 248=8+112+128
\tag{12a}
\]

The ``axial 8'' and the ``other 128'' of the four-dimensional cell count
(Eq. (6)) agree in number with the ``8 of the Cartan part'' and the
``128 half-integer roots'' of E₈. What remains on the E₈ side are the
112 integer roots. However, the 128 four-dimensional cells break down by
coordinate type into \((1,1,0,0)\) 24, \((1,1,1,0)\) 32, \((1,1,1,1)\)
16, \((2,0,0,0)\) 8 and \((2,1,0,0)\) 48, which differ in form from the
half-integer roots of E₈ (all coordinates \(\pm1/2\)). At present,
\textbf{only the numbers} agree.

\textbf{Side of evolution: 23/124 and 31, 8, 23.} The Coxeter number of
E₈ is \(h=30\), and

\[
124=4\times31=4(h+1),\qquad 248=8\times31,\qquad \frac{23\pi}{124}+\frac{8\pi}{124}=\frac{\pi}{4}
\tag{12b}
\]

hold. The angle \(23\pi/124\) of the finite-order root
\(R_{124,23}=\cos^2(23\pi/124)\) has complement \(8\pi/124\) with
respect to \(\pi/4\), and \(23+8=31=h+1\). Further, the eigenvalues of
the Coxeter element of E₈ are \(e^{2\pi i m/30}\) with
\(m\in\{1,7,11,13,17,19,23,29\}\) (the exponents of E₈), and 23 is one
of them. 8 is also the rank of E₈.

\textbf{How to read these correspondences: near but not identical.} The
correspondences of Eqs. (12a) and (12b) are all \textbf{near but not
identical}. 128 agrees in number but differs in form, and the 31 of
\(124=4\times31\) is one more than the Coxeter number 30. In the
hypothesis of this paper, this very fact of being ``near but not
identical'' is meaningful.

\begin{itemize}
\tightlist
\item
  If there were an exact correspondence between the two structures (a
  map preserving inner products and symmetry), they would lock
  completely. Then neither metastability nor sudden change would arise,
  and H2 (non-commensurability) of Chapter 4 would not hold.
\item
  If the two structures were completely unrelated, they would not pull
  each other in either. Then H3 (pull-in) would not hold.
\item
  A relation in which integers are shared while the forms disagree
  satisfies H2 and H3 at the same time. It is exactly the form required
  by the central hypothesis of this paper: pulling in but never locking
  completely.
\end{itemize}

Therefore, proving the reason why these correspondences hold (an exact
map) is not a task of this paper. If an exact map were found, it would
break the hypothesis of this paper for this case.

Under this reading, the hypothesis breaks in the following two cases.

\begin{enumerate}
\def\labelenumi{\arabic{enumi}.}
\tightlist
\item
  \textbf{Discovery of a lock:} an exact correspondence preserving inner
  products and symmetry is found between the four-dimensional 137 cells
  and the structure of E₈, or between the finite-order root 23/124 and
  the Coxeter element of E₈. The two structures are then fully
  reconciled, and this case is not an example of non-commensurability.
\item
  \textbf{Absence of pull-in:} in evolution with the two constraints
  imposed together, no bias in residence times, no phase slips and no
  metastability appear (Task 2 of Chapter 6). Closeness then has no
  effect on the dynamics, and this case is not an example of pull-in.
\end{enumerate}

\subsubsection{7.7 The Running of the Fine-Structure Constant and the
Position of
Resonances}\label{the-running-of-the-fine-structure-constant-and-the-position-of-resonances}

The low-energy \(\alpha^{-1}(0)\simeq137.036\) and the electroweak-scale
\(\alpha^{-1}(M_Z)\simeq128\) do not mean two moments of time. In
standard theory they are \textbf{effective couplings at different
momentum scales}, and the running due to vacuum polarization is
experimentally supported.

In particular, KLOE-2 measured the change of the effective
electromagnetic coupling and the hadronic contribution in the timelike
momentum region \(0.6<\sqrt s<0.975\,\mathrm{GeV}\). The appearance of
the \(\rho\)--\(\omega\) resonance is an example of interference between
different particle channels affecting the effective response. This,
however, is not a measurement showing that ``the hidden symmetries of
137 and 23/124 interfere'' as in this paper.

This case must fit \textbf{not only a numerical match at a single
low-energy point} but also the \(q\) dependence, threshold and resonance
structures, complex amplitudes and the relation to scattering cross
sections. A concrete response map for reading the observed running as
``the result of time integration'' is the central task of this case.

\subsubsection{7.8 Verification Tasks of the
Case}\label{verification-tasks-of-the-case}

\textbf{Task A: Selection theorem on the side of space.} Not only
reproduce the 137 of Eq. (5), but derive why the sector of four
dimensions, radius 3 and side 1 is allowed and other integer candidates
are not selected. If four dimensions and radius 3 remain inputs, it is
not a selection theorem.

\textbf{Task B: Selection theorem on the side of evolution.} Show that
23/124 is selected by the interaction alone from the infinitely many
rational-angle candidates of Eq. (9). Examine sensitivity to initial
values, labelling of states and arbitrary thresholds.

\textbf{Task C: Observation map to electromagnetic coupling.} Define
\(c_{\rm read}(q)\) of Task 3 in Chapter 6 as \(\alpha_{\rm read}(q)\),
and predict energy regions different from the Thomson limit with a
common parameter setting.

\textbf{Task D: Search for pairs that are near but not identical.} For
rational angles other than 23/124, exceptional structures other than E₈
and values near other physical constants, list widely the pairs
satisfying the four search conditions of Chapter 4. For each, compute
the evolution with the two constraints imposed together and examine
whether pull-in and incomplete locking occur.

\textbf{Task E: Whether the closeness to E₈ produces pull-in.} Examine
the two ways of breaking in §7.6. If an exact correspondence (a lock) is
found, this case becomes a counterexample to the hypothesis. If pull-in
and incomplete locking (bias in residence times, phase slips,
metastability) appear in evolution with the two constraints imposed
together, this case becomes an instance of the hypothesis.

If the competition between 137 and 23/124 does not hold, the closeness
of the two values remains, at least for now, a numerical coincidence.
Even then, the general hypothesis of Chapter 4 is verified independently
of this case.

\subsection{8. Similar Prior Studies and
Differences}\label{similar-prior-studies-and-differences}

\textbf{Aubry (1983) and Bak (1982).} They treat the Frenkel--Kontorova
model, commensurate and incommensurate phases, the devil's staircase and
metastable configurations. They are the most direct mathematical
precedent for the competition of different periods and the selection of
integer ratios.

\textbf{Zhu et al.~(2026).} They observed a discrete time quasicrystal
in Rydberg atoms. A modern experimental example studying several
temporal symmetries and aperiodic order. The experiment, however, uses
external driving and known time, so it does not explain the emergence of
time itself.

\textbf{Martin, Refael and Halperin (2017).} They studied time-averaged
quantized energy transfer under driving by several frequencies and in
higher-dimensional Floquet lattices. A reference example for considering
integral readout of evolution, but it takes \(\hbar\) as input and does
not seek physical constants.

The difference between these and this paper is that \textbf{this paper
tries to verify, from one unknown underlying relational structure at the
same time, several integer-selecting symmetries, their incomplete
locking, a direction of evolution without external time, and the
projection to observed constants}.

\subsection{9. Record of the Discovery Dialogue: The Order of the
Author's Shifts of
Argument}\label{record-of-the-discovery-dialogue-the-order-of-the-authors-shifts-of-argument}

The following are \textbf{short original excerpts of the author's
statements} in this dialogue, each given in English translation followed
by the original Japanese. The order of statements has not been changed
to smooth the meaning. Typos are mostly preserved in the originals.

\begin{enumerate}
\def\labelenumi{\arabic{enumi}.}
\tightlist
\item
  ``Why explain it with Planck's constant? I think the root is a
  different structure, precisely the structure that determines the
  symmetries of spacetime?'' (Original:
  「なぜプランク定数で説明するのですかね？根本は別の構造まさに時空の対称性を決める構造が根本だと思いますが？」)
\item
  ``It is complicated because several symmetries combine and require
  discreteness and periodicity in space and time; conserved quantities
  are read as particular conserved quantities by projection onto
  particular directions.'' (Original:
  「複雑なのは、複数の対称性が組み合わさって、空間　時間に離散性　と　周期性を要求するからで　保存量は特定の方向での射影で特定の保存量が読めるんですよ」)
\item
  ``A symmetry that gives magic numbers is needed; otherwise
  discretization does not arise.'' (Original:
  「マジックナンバーを与える対称性が必要です　でないと離散化が生まれません」)
\item
  ``Order and causality also need a symmetry that requires discretized
  magic numbers; without it, time and uncertainty do not arise.''
  (Original:
  「順番や因果にも離散化したマジックナンバーを要求する対称性が必要でこれがないと時間や不確定性が生まれません」)
\item
  ``At this point it is no longer the world of physical symmetries but
  of number theory.'' (Original:
  「ここまでくると物理の対称性ではなく整数論の世界です」)
\item
  ``Further, a difficult condition is needed: the several symmetries in
  fact cannot hold at the same time, and a combination is needed that
  produces sudden change and metastability in a situation where they
  pull in but never lock completely.'' (Original:
  「さらに複数の対称性は実は同時に成立できないという難しい条件が必要で、引き込むけど完全にはロックしないという状況で急激な変化と準安定を生じる組み合わせが必要です」)
\item
  ``The way I noticed this is simple: why does a complex 137 happen to
  exist near the fine-structure constant? Do you see?'' (Original:
  「これ気づいたの簡単でなぜ微細構造定数の付近に複素の137が偶然存在するのかですよ？わかりますか？」;
  ``complex'' (複素) is as in the original)
\item
  ``The second one is exactly the 23/124 of U\^{}n=I.'' (Original:
  「二つ目がまさにU \^{}n=I の23／124ですよ」)
\item
  ``So, with the 137 that determines the four-dimensional spatial
  crystal structure, the 23/124 that determines the crystallization of
  the direction of evolution, and perhaps something else, a symmetry
  constraint that may not exist, the α that oscillates metastably, read
  by time integration along the direction of evolution, becomes the
  presently observed α, I suppose.'' (Original:
  「だから4次元空間結晶構造を決める137と発展方向の結晶化を決める23/124とあと何か？ないかもしれない対称性の拘束で準安定振動するαを発展方向に時間積分して読むと現在の観測値のαになるんでしょう」)
\item
  ``This is a structure I had already noticed four months ago when I
  found the second 137; are there similar prior studies?'' (Original:
  「これは二つ目の137を見つけた4ヶ月前にすでに気づいていた構造ですが、類似性のある先行研究ありますか？」)
\item
  ``Well then, since this chat itself becomes the discovery record, let
  us write up and fix the Fourteenth Thought Experiment. What do you
  think?'' (Original:
  「んーでは第14思考実験をまさにこのチャット自身が発見記録になるので論文化して固定しておきましょう　どう考えますか？」)
\end{enumerate}

Statements 1--6 are the general principle of this paper, and statements
7--9 are what led to noticing it and the case applying the principle to
the fine-structure constant. What matters in this record is the logical
order in which the author repeatedly corrected, before the conclusion:
``not from conserved quantities'', ``do not chase \(U^n=I\)'', and
``dimension, order and causality also require integer selection''.

\textbf{Dialogue at the time of revision (2026-10-12).} The author's
statements in the revision dialogue (with Claude) before fixing v1.0 are
recorded with continuing numbers under the same policy.

\begin{enumerate}
\def\labelenumi{\arabic{enumi}.}
\setcounter{enumi}{11}
\tightlist
\item
  ``For some reason it is too fixated on α, and the paper is muddied.
  α=137 is only one case among them.'' (Original:
  「なぜかαにこだわり過ぎて、論文が濁っています。α=137はそのなかの事例の一つでしかありません。」)
\item
  ``The crystallographic restriction theorem and Hurwitz's dimension
  restriction are important, so set up a separate chapter at the end and
  add an explanation.'' (Original:
  「結晶学的制限定理やフルヴィッツの次元制限　は重要なので、末尾に別に章を立てて説明を追加しなさい。」)
\item
  ``It is rather unnatural not to take up E8, the symmetry with the most
  famous magic number. E8 is not a general symmetry requirement of
  theoretical physics, but it is a symmetry that appears in some
  reactions, isn't it?'' (Original:
  「もっとも有名なマジックナンバーのある対称性E8をとりあげていないのが逆に不自然ですよね、E８はあくまで理論物理学の一般てきな対称性の要請ではないですが、一部の反応ではあらわれる対称性ですよね？」)
\item
  ``If E8 is treated, \ldots{} it is more unnatural to avoid its
  relation to alpha?'' (Original:
  「E8を扱うと、\ldots アルファとの関係もさける方が不自然ですよ？」)
\item
  ``Σx\^{}2 = 0 is the ultimate conserved quantity or fixed point, and,
  read by central projection, there is the symmetry that the left side
  reads as observable space and the right side (or the imaginary axis on
  the left side) as the unobservable direction; further, from anonymity,
  which is also an ultimate symmetry, any axis may be the projection
  axis; further, in central projection, if the curvature is constant,
  the origin is shared; this symmetry too should be explained here. In a
  sense, in my theoretical system, this is the ultimate symmetry.''
  (Original: 「Σx\^{}2 = 0
  って、究極の保存量または、固定点であるとともに、中心投影の読み方をすると、左辺が観測可能空間で、右辺（または左辺の虚数軸）が観測できない方向によめるという、対称性と、さらには、無名性というこれも究竟の対称性から、どの軸を、投影軸にしてもよい、さらに、中心投影では、曲率が一定なら、原点が共有されるという対称性もここで、説明しておくべきですよ、ある意味私の理論体系では、これが究極の対称性ですよ」)
\item
  ``Σxn\^{}2 = 0, like U\^{}n=I, is rather not an assumption but derived
  as a result of combining other symmetries \ldots{} In fact, in past
  experiments too, without placing Σxn\^{}2=0 or U\^{}n=I as assumptions
  to derive from, they turned out to hold as a result; this is an
  empirical rule.'' (Original: 「Σxn\^{}2 = 0 も、U\^{}n=I
  と同様に、むしろ仮定ではなく、他の対称性の組み合わせで結果的に導出される\ldots 実際も過去の実験でも、Σxn\^{}2=0
  、U\^{}n=I
  を仮定として置いて導出しなくても、結果的に、それが成り立っていたという経験則です。」)
\item
  ``It is not that there is a reason for it to hold; being in the
  vicinity is what is meaningful, and I am hypothesizing that it does
  not lock completely but causes pull-in, producing rapid development
  and metastability; so, conversely, if there were a mechanism that
  locks, this hypothesis would collapse.'' (Original:
  「成立する理由があるのではなく、近傍になることにいみがあり、それが、完全にはロックしないが、引き込みを起こして、急激な発展や、準安定を生むのではないか、と仮定しているので、逆にロックする仕組があったら、この仮説が崩れますよ」)
\item
  ``It is a proposal: why not search for such relations, seemingly not
  theoretical, strangely unrelated yet close, causing pull-in but never
  locking completely?'' (Original:
  「この一見理論的でない、不思議な無関係だけど、近くにあり、引き込みは起こすけど、完全ロックしない、という関係をさがしてはどうか？という提案なんですよ」)
\item
  ``Σxn\^{}2 = 0 can be said as 0 + 0 = 0, so in fact anything looks
  like a conserved quantity, that is what I am saying. \ldots{} If
  Σxn\^{}2 = R\^{}2, then Σxn\^{}2 + (iR)\^{}2 = 0, and as a result
  Σx'\^{}2 = 0, right? That is, it does not give a conserved quantity of
  0; it is only a definition saying that a conserved quantity can
  exist.'' (Original: 「Σxn\^{}2 = 0 って、 0 + 0 = 0
  といっちえるので、実は何でも保存量にみえるといっています。\ldots Σxn-2
  = R\^{}2 とすると、　Σxn\^{}2 + (iR)\^{}2 = 0 となり、結果敵に、
  Σx'\^{}2 = 0
  ですよね　つまり保存料０をあたえるのではなく、保存量が存在し得るといっているだけの定義なんですよ？」)
\end{enumerate}

Statements 12--20 determined the content of the paper in the revision:
that α is one case (statement 12); precedents for integer selection and
E₈ (statements 13--15); that the zero square sum is the ultimate
symmetry, a consequence rather than an assumption, and a definition that
a conserved quantity can exist rather than the value of a conserved
quantity (statements 16, 17, 20); and that being near but not identical
is the content of the hypothesis, with the proposal to search for such
pairs being the core of this paper's claim (statements 18, 19).

\subsection{10. Precedents in Which the Compatibility of Two Structures
Selects Integers: The Crystallographic Restriction Theorem, Hurwitz's
Theorem and
E₈}\label{precedents-in-which-the-compatibility-of-two-structures-selects-integers-the-crystallographic-restriction-theorem-hurwitzs-theorem-and-eux2088}

The hypothesis of this paper (§3.2--3.4) rests on the view that integers
are selected not by a single symmetry but \textbf{only when several
structures of different character are imposed at the same time}.
Mathematics has at least three rigorous precedents for this view: the
crystallographic restriction theorem, Hurwitz's theorem, and E₈. E₈ is
also an example in which the selected structure appears in real matter
only under particular conditions. This chapter describes for each the
content, the selected integers, and the correspondence to and limits
with respect to the hypothesis of this paper.

\subsubsection{10.1 The Crystallographic Restriction Theorem:
Compatibility of Rotation and
Translation}\label{the-crystallographic-restriction-theorem-compatibility-of-rotation-and-translation}

\textbf{Content.} The rotations that preserve a lattice in the plane or
in three-dimensional space (a periodic structure with translational
symmetry) are limited to those of order 1, 2, 3, 4, 6. Five-fold
rotations and rotations of order seven or more are compatible with no
lattice.

\textbf{Reason.} Written in the basis of the lattice, the matrix of a
rotation has all integer entries. Hence its trace is an integer. On the
other hand, if the rotation angle is \(\theta\), the trace is
\(2\cos\theta\) in the plane and \(1+2\cos\theta\) in three dimensions.
In either case \(2\cos\theta\) must be an integer, so

\[
2\cos\theta\in\{-2,-1,0,1,2\}
\quad\Longrightarrow\quad
\theta\in\left\{0,\ \tfrac{\pi}{3},\ \tfrac{\pi}{2},\ \tfrac{2\pi}{3},\ \pi\right\}
\tag{13}
\]

If the rotation angle is \(2\pi/m\), the allowed orders are limited to
\(m=1,2,3,4,6\).

\textbf{Selected integers.} Rotation (the discrete part of a continuous
symmetry) and translation (the periodicity of the lattice) each forbid
no order by themselves. The integers \(\{1,2,3,4,6\}\) appear
\textbf{only when both are required at the same time}.

\textbf{Extension to dimensions.} In a general dimension \(d\) as well,
a finite-order rotation preserving a lattice needs a minimal dimension
depending on its order \(m\). Let \(\psi(m)\) be the function given by
the sum of Euler's function over the prime-power factors of \(m\)
(subtracting 1 when \(m\equiv2\pmod 4\)); a rotation of order \(m\) is
compatible with a \(d\)-dimensional lattice only when \(\psi(m)\le d\).
In four dimensions the allowed orders widen to \(1,2,3,4,5,6,8,10,12\).
\textbf{Which integers are allowed changes with the dimension of space.}

\textbf{When incompatible: quasicrystals.} Five-fold symmetry and
periodic translation are incompatible in three dimensions. Yet in 1984
an alloy showing a diffraction pattern with five-fold (icosahedral)
symmetry was observed. This is a structure (a quasicrystal) that has
lost periodicity but keeps long-range order, and mathematically it can
be described as a structure obtained by projecting a higher-dimensional
lattice to a lower dimension. \textbf{When incompatible symmetries are
required at the same time, the result is not disorder but another order
that can be read as a projection from higher dimensions.}

\paragraph{\texorpdfstring{10.1.1 Application to Phase: Only \(Z_4\) and
\(Z_6\) Can Be Marked Exactly with
Integers}{10.1.1 Application to Phase: Only Z\_4 and Z\_6 Can Be Marked Exactly with Integers}}\label{application-to-phase-only-z_4-and-z_6-can-be-marked-exactly-with-integers}

Represent the phase of an oscillation by a pair of integers \((a,b)\) (a
point of the two-dimensional integer lattice, for example \(a+bi\)).
Suppose the operation of rotating this phase by a fixed angle is to be
implemented exactly by integer addition, subtraction and exchange only.
That operation is a \(2\times2\) matrix with integer entries and
determinant \(\pm1\), and, for the same reason as Eq. (13), the orders
of those that return after finitely many steps are limited to 1, 2, 3,
4, 6. The implementations fall into two lattices.

{\def\LTcaptype{none} % do not increment counter
\begin{longtable}[]{@{}lllrr@{}}
\toprule\noalign{}
Lattice & Step & Integer map & Trace & Order \\
\midrule\noalign{}
\endhead
\bottomrule\noalign{}
\endlastfoot
\(Z_4\) (Gaussian integers) & 90° & \((a,b)\mapsto(-b,a)\) & 0 & 4 \\
\(Z_6\) (Eisenstein integers) & 60° & \((a,b)\mapsto(a-b,a)\) & 1 & 6 \\
\end{longtable}
}

The \(Z_4\) map consists only of exchanging the components and reversing
the sign of one of them, and corresponds to multiplication by \(i\).
Applying the \(Z_6\) map twice gives 120° (\(Z_3\)). With the attached
\texttt{verify\_\allowbreak{}phase\_\allowbreak{}lattice.py} it was confirmed that, for all
integer matrices with determinant \(\pm1\) and entries of absolute value
at most 4, the finite orders are only \(\{1,2,3,4,6\}\).

45° (order 8) is the first angle that cannot be marked exactly on the
two-dimensional integer lattice (it is possible on a four-dimensional
lattice, corresponding to \(\psi(8)=4\) in §10.1). In quantum
computation, circuits made of Hadamard gates, CNOT and 90° phase gates
(Clifford circuits) can be tracked efficiently by classical computers
(the Gottesman--Knill theorem {[}13{]}), and adding the 45° phase gate
(the T gate) makes them universal. The phase steps of Clifford circuits
agree with the 90° of \(Z_4\), and the first angle outside them is 45°.
The boundary of phases that can be marked exactly with integers and the
boundary of circuits that can be tracked classically are drawn at the
same angle. This paper records this correspondence as a concrete clue
for integer selection on the side of evolution (§3.3(c)).

\subsubsection{10.2 Hurwitz's Theorem: Compatibility of Multiplication
and
Norm}\label{hurwitzs-theorem-compatibility-of-multiplication-and-norm}

\textbf{Content.} The algebras over the reals whose norm preserves
products (normed division algebras satisfying \(|xy|=|x||y|\), without
zero divisors) are limited to dimensions 1, 2, 4, 8: the real numbers,
complex numbers, quaternions and octonions.

\textbf{Equivalent form: sum-of-squares identities.} This is the same as
saying that the \(n\) for which

\[
(x_1^2+\cdots+x_n^2)(y_1^2+\cdots+y_n^2)=z_1^2+\cdots+z_n^2
\tag{14}
\]

can be satisfied identically with each \(z_k\) bilinear in \(x\) and
\(y\) are limited to \(1,2,4,8\). \(n=2\) corresponds to the absolute
value of complex numbers, \(n=4\) to Euler's four-square identity, and
\(n=8\) to Degen's eight-square identity.

\textbf{Selected integers.} Multiplication alone (algebraic structure)
or norm alone (metric structure) can be made in any dimension. The
integers \(\{1,2,4,8\}\) appear \textbf{only when product and norm are
made compatible at the same time}.

\textbf{Properties lost as the dimension is raised.} Each time the
dimension is doubled, one property is lost.

{\def\LTcaptype{none} % do not increment counter
\begin{longtable}[]{@{}rll@{}}
\toprule\noalign{}
Dimension & Algebra & Property newly lost \\
\midrule\noalign{}
\endhead
\bottomrule\noalign{}
\endlastfoot
1 & real numbers & ― \\
2 & complex numbers & ordering \\
4 & quaternions & commutativity \\
8 & octonions & associativity \\
16 & sedenions & multiplicativity of the norm (zero divisors appear) \\
\end{longtable}
}

In 16 dimensions, the compatibility of product and norm itself fails.

\textbf{Related results.} The same integers \(1,2,4,8\) appear in the
parallelizability of the spheres \(S^{n-1}\) (only \(S^0,S^1,S^3,S^7\))
and in the existence of maps of Hopf invariant one (Adams, 1960). That
mutually different questions select the same set of integers shows that
the integers stem not from the individual questions but from the
compatibility conditions of structures.

\subsubsection{10.3 E₈: The Compatibility of Evenness and Self-Duality
Selects 8, and a Symmetry That Appears in Particular
Reactions}\label{eux2088-the-compatibility-of-evenness-and-self-duality-selects-8-and-a-symmetry-that-appears-in-particular-reactions}

E₈ is the best-known ``symmetry with a magic number''. E₈ is not a
symmetry that theoretical physics requires in general, nor is it
contained in the gauge group of the Standard Model. Yet in mathematics
it is selected uniquely by the compatibility of two structures, and in
physics it appears only under particular conditions. A single object
shows both sides of the theme of this paper.

\textbf{Mathematics: the compatibility of evenness and self-duality
selects dimension 8.} Impose two conditions on a lattice.

\begin{itemize}
\tightlist
\item
  \textbf{Even:} the squared length of every lattice vector is an even
  integer.
\item
  \textbf{Self-dual (unimodular):} the lattice coincides with its own
  dual lattice. The volume of the fundamental domain is 1.
\end{itemize}

Positive-definite lattices satisfying both at the same time exist only
in dimensions that are multiples of 8. In 8 dimensions there is exactly
one, the E₈ lattice (in 16 dimensions there are 2, in 24 dimensions the
24 Niemeier lattices). Each condition alone forbids no dimension, but
making them compatible selects 8. This is a selection of the same form
as in §10.1 and §10.2. The integers that the two integer structures of
the case in Chapter 7 share with E₈ are recorded in §7.6.

The shortest vectors (roots) of the E₈ lattice number 240 and split into
two types:

\[
240=112+128
\tag{15}
\]

The 112 are of the type with two coordinates \(\pm1\) and the rest 0;
the 128 are of the type with all coordinates \(\pm1/2\) and an even
number of minus signs. Adding eight mutually commuting directions (the
Cartan part) gives the dimension \(248=8+240\) of the Lie group E₈. In
sphere packing in 8 dimensions the E₈ lattice is the densest (Viazovska
2017 {[}15{]}), and the number of spheres touching one sphere is 240.

\textbf{Connection with Hurwitz's theorem.} Within the octonions, which
correspond to the last dimension 8 of §10.2, the set of elements
corresponding to ``integers'' (the integral octonions) forms the E₈
lattice (Coxeter 1946 {[}14{]}). The 8 selected by the compatibility of
product and norm and the 8 selected by the compatibility of evenness and
self-duality meet in the same lattice.

\textbf{Connection with quasicrystals.} Projecting the root system of E₈
suitably to four dimensions gives a four-dimensional quasicrystal with
the symmetry \(H_4\) of the 600-cell (Elser and Sloane 1987 {[}16{]}).
As seen in §10.1, five-fold symmetry is incompatible with
three-dimensional periodic lattices. That five-fold symmetry appears,
with order preserved, as a projection from the eight-dimensional E₈. The
golden ratio \((1+\sqrt5)/2=2\cos(\pi/5)\) appears in the construction
of \(H_4\).

\textbf{Physics: E₈ appearing only in particular reactions.} The
one-dimensional transverse-field Ising model is at a quantum critical
point when the transverse field takes a certain value. Adding a weak
longitudinal field there turns the low-energy excitations into eight
kinds of particles, whose mass ratios are determined by the structure of
E₈ (Zamolodchikov 1989 {[}17{]}). The ratio of the two lightest masses
is

\[
\frac{m_2}{m_1}=2\cos\frac{\pi}{5}=1.618\ldots
\tag{16}
\]

equal to the golden ratio. In 2010, neutron scattering on the
quasi-one-dimensional Ising magnet CoNb₂O₆ observed that the mass ratio
of the two lightest excitations near the critical point approaches the
golden ratio (Coldea et al.~2010 {[}18{]}). Since then, E₈ spectra have
also been reported in BaCo₂V₂O₈ (Zou et al.~2021 {[}19{]}).

This E₈ is not a symmetry of the fundamental laws that make matter. It
appears in the low-energy excitations only at \textbf{a particular point
combining two conditions}: the critical point due to the transverse
field, and the longitudinal field that moves away from it. This fits the
view of this paper that integer structures are selected from
combinations of conditions, rather than a single symmetry being given at
the root. It is also recorded that the golden ratio that appeared in
quasicrystals appears here as a mass ratio.

\textbf{E₈ as a consistency condition of theory.} In superstring theory,
the condition of anomaly cancellation selects \(SO(32)\) and
\(E_8\times E_8\) as gauge groups (Green and Schwarz 1984 {[}20{]}). The
heterotic string with \(E_8\times E_8\) was constructed from this
condition (Gross et al.~1985 {[}21{]}). Here too E₈ is selected as a
consistency condition of the framework of placing strings in a
background spacetime.

\subsubsection{10.4 Correspondence to and Limits with Respect to the
Hypothesis of This
Paper}\label{correspondence-to-and-limits-with-respect-to-the-hypothesis-of-this-paper}

{\def\LTcaptype{none} % do not increment counter
\begin{longtable}[]{@{}
  >{\raggedright\arraybackslash}p{(\linewidth - 8\tabcolsep) * \real{0.2000}}
  >{\raggedright\arraybackslash}p{(\linewidth - 8\tabcolsep) * \real{0.2000}}
  >{\raggedright\arraybackslash}p{(\linewidth - 8\tabcolsep) * \real{0.2000}}
  >{\raggedright\arraybackslash}p{(\linewidth - 8\tabcolsep) * \real{0.2000}}
  >{\raggedright\arraybackslash}p{(\linewidth - 8\tabcolsep) * \real{0.2000}}@{}}
\toprule\noalign{}
\begin{minipage}[b]{\linewidth}\raggedright
Aspect
\end{minipage} & \begin{minipage}[b]{\linewidth}\raggedright
Crystallographic restriction theorem
\end{minipage} & \begin{minipage}[b]{\linewidth}\raggedright
Hurwitz's theorem
\end{minipage} & \begin{minipage}[b]{\linewidth}\raggedright
E₈
\end{minipage} & \begin{minipage}[b]{\linewidth}\raggedright
Hypothesis of this paper
\end{minipage} \\
\midrule\noalign{}
\endhead
\bottomrule\noalign{}
\endlastfoot
Structures imposed together & rotation and translation & multiplication
and norm & evenness and self-duality & constraints on the side of space
and on the side of evolution \\
Selected integers & \(\{1,2,3,4,6\}\) (depending on dimension) &
\(\{1,2,4,8\}\) & dimension 8 (multiples of 8), 240 roots & unspecified
(in the case of Chapter 7, candidates 137 and 23/124) \\
When incompatible & quasicrystal (order readable as projection from
higher dimensions) & multiplicativity of the norm breaks in 16
dimensions & does not exist in dimensions other than multiples of 8 &
metastability that pulls in but never locks completely (H2--H4) \\
Appearance in physics & crystals, quasicrystals & ― & excitation
spectrum under the particular condition of critical point plus
longitudinal field & readout of metastable oscillation \\
Time and evolution & static compatibility condition & static
compatibility condition & static compatibility condition (the physical
example is the spectrum near the critical point) & direction of
evolution without external time (Chapter 5) \\
\end{longtable}
}

The three precedents give a rigorous mathematical basis to the starting
point of this paper (§3.2--3.4) that \textbf{integers are selected not
by a single symmetry but by the compatibility conditions of several
structures}. In particular, the crystallographic restriction theorem and
E₈ show that, even when incompatible symmetries are required at the same
time, order remains and can be read as a projection from higher
dimensions. This is the precedent closest to the ``projective readout''
and ``incomplete locking'' of this paper. E₈ further shows that the
structure so selected appears in real matter only under a particular
combination of conditions.

Applying the crystallographic restriction to both space and phase gives
one concrete example of mismatch. The rotation orders allowed by
four-dimensional spatial lattices are \(\{1,2,3,4,5,6,8,10,12\}\), and
the orders allowed for phases represented by a pair of integers are
\(\{1,2,3,4,6\}\); the two do not agree. On the side of space and on the
side of evolution, the same theorem selects different sets of integers.
Whether this mismatch becomes a collision of operations on the same
state is Task 4 of Chapter 6.

On the other hand, these precedents deal with static compatibility
conditions and do not contain \textbf{non-commensurability within
evolution}, such as pull-in, phase slips and metastability. The
hypothesis of this paper asks whether these precedents can be extended
to dynamics along the direction of evolution without external time.

\subsection{\texorpdfstring{11. The Ultimate Symmetry: The Zero Square
Sum \(\sum x^2=0\) and Central
Projection}{11. The Ultimate Symmetry: The Zero Square Sum \textbackslash sum x\^{}2=0 and Central Projection}}\label{the-ultimate-symmetry-the-zero-square-sum-sum-x20-and-central-projection}

The symmetries treated in this paper (rotation of lattices,
multiplication and norm, evenness and self-duality, constraints on the
side of space and of evolution) were all integer-selecting symmetries.
In the author's theoretical system, one symmetry is placed above them.
It is the zero square sum

\[
\sum_{i=1}^{n}x_i^2=0
\tag{17}
\]

This is the same equation introduced as the equation of the light cone
in the Twelfth Thought Experiment {[}K6{]}. This chapter first states
that this equation holds not as an assumption but as a consequence of
combining other symmetries (§11.1), and then that this single equation
has four symmetries at the same time.

\subsubsection{11.1 Holding Not as an Assumption but as a
Consequence}\label{holding-not-as-an-assumption-but-as-a-consequence}

Eq. (17) is not an axiom placed at the starting point of this paper.
Like \(U^n=I\) (item 2 of §0), \textbf{it holds as a result of combining
other symmetries}. In the author's numerical experiments so far, too,
neither \(U^n=I\) nor \(\sum x^2=0\) was placed as an input. It was
confirmed afterwards that they held as consequences of the interaction
and integer conditions (for \(U^n=I\), the Thirteenth Thought Experiment
{[}K7{]} confirmed that the period was a consequence of the input
\(R\)). This is an empirical rule, and mathematics has two rigorous
facts of the same form.

\textbf{\(U^n=I\) comes from the combination of integrality and norm
preservation.} Transformations that preserve a lattice (can be written
with integer entries) and transformations that preserve the norm
(orthogonal or unitary) do not, each by themselves, require returning
after finitely many steps. But the set of transformations satisfying
both at the same time is a finite group, and for every element \(U\) of
it there is a finite \(n\) with

\[
U^n=I
\tag{17a}
\]

In three dimensions there are 48 matrices with integer entries that
preserve length, and their orders are limited to 1, 2, 3, 4, 6. This is
the same set as the crystallographic restriction of §10.1. Finite
recurrence is not an assumption but a consequence of imposing two
symmetries at the same time.

\textbf{\(\sum x^2=0\) comes from the compatibility of multiplication
and norm.} Squaring a complex number \(Z=a+ib\),

\[
(\mathrm{Re}\,Z^2)^2+(\mathrm{Im}\,Z^2)^2=|Z|^4,
\qquad
(a^2-b^2)^2+(2ab)^2=(a^2+b^2)^2
\tag{17b}
\]

holds identically, regardless of \(a,b\). Placing \(|Z|^2\) on the
imaginary axis, \((a^2-b^2,\ 2ab,\ i(a^2+b^2))\) is a solution of Eq.
(17). Eq. (17) is not placed anywhere as an assumption. It holds
automatically from the fact that the norm of a product equals the
product of the norms (\(|Z^2|=|Z|^2\), the compatibility of
multiplication and norm in §10.2). Integer \((a,b)\) give all primitive
Pythagorean triples (Euclid's formula; \((a,b)=(2,1)\) gives \(3,4,5\)).

The same construction extends to pairs of two components. From a pair
\((p,q)\) of elements of the algebras of dimension \(k=1,2,4,8\) allowed
by Hurwitz's theorem (real numbers, complex numbers, quaternions,
octonions), form

\[
t=|p|^2+|q|^2,\qquad z=|p|^2-|q|^2,\qquad w=2\bar p\,q
\tag{17c}
\]

Then \(t^2-z^2-|w|^2=4\left(|p|^2|q|^2-|\bar p\,q|^2\right)=0\) holds
identically, and points on the light cone of dimension \(k+2=3,4,6,10\)
are obtained {[}22{]}{[}23{]}. The case \(k=2\) is the correspondence
between the four-dimensional light cone and spinors. In the algebra of
dimension 16 (the sedenions) the norm does not preserve products, so
this identity fails (confirmed by the attached re-verification).
\textbf{The light cone appears without assumption only in the dimensions
in which multiplication and norm are compatible.}

\textbf{The meaning of ``ultimate''.} Eq. (17) is called the ultimate
symmetry not because it is the starting point (axiom) of everything
else, but because it is \textbf{the common form} that appears as a
result when structures of different character, such as integrality, norm
and multiplication, are combined. That form does not fix the value of a
conserved quantity; it expresses only that a conserved quantity can
exist (§11.3). \(U^n=I\) and \(\sum x^2=0\) are both ``relations that
hold as results of combination'', and placing them at the starting point
reverses cause and effect. The same discipline by which this paper, in
item 2 of Chapter 0, does not give \(U^n=I\) in advance as an
explanatory principle is applied to \(\sum x^2=0\) as well.

\subsubsection{11.2 Nontrivial Solutions Need an Imaginary
Axis}\label{nontrivial-solutions-need-an-imaginary-axis}

Within the real numbers, the only solution of Eq. (17) is \(x=0\). For a
nonzero solution to exist, the square of at least one axis must be
negative, that is, its component must be imaginary. And one imaginary
axis suffices. Setting one imaginary axis \(x_n=it\), Eq. (17) becomes

\[
\sum_{i=1}^{n-1}x_i^2=t^2
\tag{18}
\]

which is the equation of the light cone itself. Its integer solutions
are Pythagorean tuples such as \(3^2+4^2=5^2\); the structure of
integers is in this equation from the outset.

\subsubsection{11.3 A Definition That a Conserved Quantity Can Exist (0
+ 0 = 0)}\label{a-definition-that-a-conserved-quantity-can-exist-0-0-0}

The 0 of Eq. (17) is not the value of a conserved quantity. A system
whose conserved quantity has the value \(R^2\),

\[
\sum_i x_i^2=R^2
\]

can be rewritten, by adding the value \(R\) as one axis \(iR\), as

\[
\sum_i x_i^2+(iR)^2=0,\qquad\text{that is}\qquad \sum_j x_j'^2=0
\tag{18a}
\]

A conserved quantity of any value takes the form of Eq. (17) when one
axis is added. Eq. (17) therefore does not fix the value of a conserved
quantity to 0; it is a definition expressing only that \textbf{a
conserved quantity can exist} (that the system is closed). The value has
not disappeared; it has moved onto the added axis.

Since \(0+0=0\), combining systems that satisfy Eq. (17) gives a
combined system that again satisfies Eq. (17). However many closed
systems are combined, the result stays closed, and in this sense
anything looks like a conserved quantity. 0 itself carries no
information. The information lies in which axis carries the value
(\(R\)).

Multiplying everything by \(\lambda\) does not change Eq. (17). Since
the value \(R\) has become one of the axes, Eq. (17) needs no external
scale. Scale is read as an internal ratio of the system.

\subsubsection{11.4 The Reading by Central Projection: Observable and
Unobservable
Directions}\label{the-reading-by-central-projection-observable-and-unobservable-directions}

Regard the nonzero solutions of Eq. (18) as a collection of half-lines
(rays) drawn from the origin. Under central projection from the origin,
\(x_i\mapsto x_i/t\), each ray is mapped to

\[
\sum_{i=1}^{n-1}\left(\frac{x_i}{t}\right)^2=1
\]

that is, to one point of the unit sphere of dimension \(n-2\). The
position on the sphere (the direction) can still be read after
projection. The position along the ray (the size of \(t\)) is folded
into a single point by the projection and can no longer be read. Integer
solutions (Pythagorean tuples) are mapped by the projection to rational
points on the sphere.

In this reading, the left side of Eq. (18) corresponds to
\textbf{observable space}, and the right side (the imaginary axis when
moved to the left side) to \textbf{the unobservable direction}. The very
division into observable and unobservable directions is built into Eq.
(17). The loss of one dimension under projection is not an approximation
added from outside but a consequence of reading Eq. (17) by central
projection. It is by this reading that Chapter 5 set the observed
constants as a projective readout of the history along the direction of
evolution.

\subsubsection{11.5 Anonymity: Any Axis May Be the Projection
Axis}\label{anonymity-any-axis-may-be-the-projection-axis}

Viewed over the complex numbers, all axes of Eq. (17) are equivalent.
Permutations of axes, and indeed all complex orthogonal transformations,
preserve Eq. (17). The equation itself does not decide which axis is
chosen as the imaginary axis (the projection axis). Solutions obtained
with different choices are carried into one another by permutations of
axes.

That the axes have no names (labels), that is, \textbf{anonymity}, is
the most fundamental symmetry in the author's theoretical system. Eq.
(17) has this anonymity as it is. That a certain axis is read as
``time'' is the result of choosing that axis as the projection axis, not
because the axis bears the name of time from the outset.

\subsubsection{11.6 If the Curvature Is Constant, the Origin Is
Shared}\label{if-the-curvature-is-constant-the-origin-is-shared}

On a surface of constant curvature, the origin of central projection is
shared by all points. On a sphere, the normals at all points pass
through one point (the center). Conversely, a connected surface all of
whose normals pass through one point is limited to part of a sphere
centered at that point (since motion perpendicular to the normal does
not change the distance from that point).

The same holds for constant negative curvature. With the right side of
Eq. (18) set to a constant,

\[
t^2-\sum_{i=1}^{n-1}x_i^2=R^2
\tag{19}
\]

is a hyperbolic space of curvature \(-1/R^2\), and central projection
from the origin uses the same origin at all points (the Beltrami--Klein
model). As in §11.3, adding the value \(R\) as one axis turns Eq. (19)
into \(\sum_i x_i^2+(it)^2+R^2=0\), that is, a zero square sum of one
higher dimension. The sphere \(\sum_i x_i^2=R^2\) likewise becomes
\(\sum_i x_i^2+(iR)^2=0\). Every surface of constant curvature is a
section of a zero square sum of one higher dimension and shares the
origin of that zero square sum.

If the curvature is constant, an observer standing at any point sees the
projection from the same origin, and there is no special observer. The
sharing of the origin expresses geometrically the anonymity of observers
(that all points are equivalent). Eq. (17) represents these surfaces
that share the origin together, in one form in which the value \(R\) has
been moved to an axis.

\subsubsection{11.7 Relation to the Hypothesis of This
Paper}\label{relation-to-the-hypothesis-of-this-paper}

None of the four symmetries (that a conserved quantity can exist, the
division into observable and unobservable directions, the anonymity of
axes, and the sharing of the origin) selects a specific integer. What
Eq. (17) selects is a \textbf{framework}: the minimal number of
imaginary axes (one) and the shape of the light cone, not the values of
\(\alpha\), masses or the number of generations.

Eq. (17) does not fix values. It only gives the framework in which a
conserved quantity can exist. What holds exactly is this layer of the
framework; on top of it, the layer that selects integers is near but not
identical and never locks completely. The statement in item 3 of Chapter
0 that ``searching only for a common fixed point is not enough''
concerns this layer of integer selection.

The symmetries of Chapters 3--10 are integer-selecting symmetries on top
of this framework. In the author's outlook, the structure is layered in
the following order.

\begin{enumerate}
\def\labelenumi{\arabic{enumi}.}
\tightlist
\item
  \textbf{Eq. (17):} the ultimate symmetry, with neither scale nor
  names.
\item
  \textbf{Central projection:} choosing one projection axis separates
  observable space from the unobservable direction.
\item
  \textbf{Integer-selecting symmetries:} integers are selected on the
  observable side (spatial lattices) and on the unobservable side
  (phases of evolution) respectively (§3.3, §10.1.1).
\item
  \textbf{Non-commensurability and pull-in:} the integer selections of
  the two sides cannot be fully reconciled, and metastable evolution
  arises (Chapter 4).
\item
  \textbf{Readout:} the history read from the direction of projection
  becomes the observed constants (Chapter 5).
\end{enumerate}

\(\mathcal C_{\rm A}\) (the side of space) and \(\mathcal C_{\rm B}\)
(the side of evolution) of Chapter 4 correspond to the observable and
unobservable sides when Eq. (17) is read by central projection. That the
two constraints cannot be fully reconciled may be a consequence of
reading one and the same equation from two sides. Showing this
possibility is the deepest candidate this paper has for Task 1 of
Chapter 6.

\subsection{12. Conclusion}\label{conclusion}

\begin{enumerate}
\def\labelenumi{\arabic{enumi}.}
\tightlist
\item
  The question of this paper is \textbf{why symmetries that select
  different integer conditions are required at the same time, why they
  cannot be fully reconciled, and why they are nevertheless pulled in by
  interaction}.
\item
  The proposal of this paper is to carry out this question as \textbf{a
  search strategy}. Search for pairs that appear theoretically unrelated
  yet lie close, and that pull in but never lock completely. The reason
  for the closeness is not explained. A pair for which an exact
  agreement is found is removed from the candidates.
\item
  If that incomplete locking produces metastable oscillation and sudden
  change, and the integral readout along the direction of evolution
  connects to the observed constants, this becomes a concrete candidate
  for thinking about the origin of quantum structure without first
  assuming Planck's constant.
\item
  Whether symmetries that require integers for order and causality are
  the origin of time and uncertainty is included in the question of this
  paper.
\item
  The 137-cell count of the four-dimensional lattice and the 23/124
  finite recurrence in the two-wave exchange system, treated as a case,
  are separate mathematical results that can be recalculated. The
  numerical closeness to the fine-structure constant was what prompted
  the idea, not evidence that both explain the same electromagnetic
  phenomenon.
\item
  From here on, verification proceeds in the order: proof of
  non-commensurability, computation of metastable dynamics, the
  observation map, integer selection of order and causality, and the
  selection theorems for each case together with comparison with the
  standard running.
\item
  Above these integer-selecting symmetries, the zero square sum
  \(\sum x^2=0\) is placed as the ultimate symmetry. It is, however, not
  an axiom but, like \(U^n=I\), a common form holding as a consequence
  of combining integrality, norm and multiplication. This equation is a
  definition expressing that a conserved quantity can exist, rather than
  fixing the value of a conserved quantity to 0; it separates observable
  and unobservable directions under central projection, any axis may be
  the projection axis (anonymity), and under constant curvature it
  shares the origin (Chapter 11).
\end{enumerate}

\textbf{The shortest statement of this thought experiment:}

\begin{quote}
Several symmetries that select special integers cannot be fully
reconciled and continue a metastable evolution while pulling each other
in. Is not what is read, by integrating the history of that incomplete
locking from a particular direction, the observed physical constants?
The fine-structure constant is one example. What should therefore be
searched for are pairs that appear theoretically unrelated yet lie
close, and that pull in but never lock completely.
\end{quote}

\begin{center}\rule{0.5\linewidth}{0.5pt}\end{center}

\subsection{References and Prior
Materials}\label{references-and-prior-materials}

\subsubsection{The Author's Earlier Work and Discovery
Materials}\label{the-authors-earlier-work-and-discovery-materials}

\begin{itemize}
\tightlist
\item
  {[}K1{]} Noriaki Kihara, \texttt{project\_\allowbreak{}w10\_\allowbreak{}alpha137.md},
  2026-04-21 to 05-30, research record on the three-dimensional
  representation of 137 and the four-dimensional cell count. Stored in
  Google Drive.
\item
  {[}K2{]} Noriaki Kihara, \texttt{思考実験\_\allowbreak{}αのエネルギー依存性.md}
  (thought experiment: energy dependence of α), 2026-05-14, exploratory
  examination of the 137 ground structure and the running.
  \href{https://drive.google.com/file/d/1yhH5Wk-Qi2IojaEgZLc_P2V2Z9Cng1nH/view}{Google
  Drive}.
\item
  {[}K3{]} Noriaki Kihara,
  \texttt{20260717\_\allowbreak{}axiomatic\_\allowbreak{}approach\_\allowbreak{}v2.md}, 2026-07-17, the
  four-dimensional cell count 137, the 128 shell, the numerical readout
  axiom and independent calculation.
  \href{https://drive.google.com/file/d/1FGNsJHZqMyTWUm6pwqWxDbpfq3BvZa6T/view}{Google
  Drive}.
\item
  {[}K4{]} Noriaki Kihara,
  \texttt{反復交換散乱における有限位数共鳴の発見...} (discovery of
  finite-order resonance in iterated exchange scattering), analysis and
  numerical results of finite order.
  \href{https://drive.google.com/file/d/1F19s-4mlGYB93IkSYrNtS9CUfGvaOFs3/view}{Google
  Drive}.
\item
  {[}K5{]} Noriaki Kihara,
  \texttt{project\_\allowbreak{}alpha\_\allowbreak{}transmission\_\allowbreak{}charge\_\allowbreak{}bridge.md}, from
  2026-07-30, research record examining the relation between two values
  near \(\alpha\) and finite-order roots.
  \href{https://drive.google.com/file/d/1F2RUrAS7Aj6xTwU9MbVm-JfqW5dgK6yo/view}{Google
  Drive}.
\item
  {[}K6{]} Noriaki Kihara, The Twelfth Thought Experiment, 2026-10-08,
  DOI:
  \href{https://doi.org/10.5281/zenodo.23226258}{10.5281/zenodo.23226258}.
\item
  {[}K7{]} Noriaki Kihara, The Thirteenth Thought Experiment v2.0,
  2026-10-10, DOI:
  \href{https://doi.org/10.5281/zenodo.23280070}{10.5281/zenodo.23280070};
  \href{https://drive.google.com/file/d/10SJBaaaI7OktbwF6-0Q_PYHR8v_XcJr8/view}{Markdown
  manuscript}.
\end{itemize}

\subsubsection{External Studies}\label{external-studies}

\begin{itemize}
\tightlist
\item
  {[}1{]} P. Bak, ``Commensurate phases, incommensurate phases and the
  devil's staircase,'' \emph{Reports on Progress in Physics}
  \textbf{45}, 587 (1982). DOI:
  \href{https://doi.org/10.1088/0034-4885/45/6/001}{10.1088/0034-4885/45/6/001}.
\item
  {[}2{]} S. Aubry, ``The twist map, the extended Frenkel-Kontorova
  model and the devil's staircase,'' \emph{Physica D} \textbf{7},
  240--258 (1983). DOI:
  \href{https://doi.org/10.1016/0167-2789(83)90129-X}{10.1016/0167-2789(83)90129-X}.
\item
  {[}3{]} D.-Y. Zhu et al., ``Observation of discrete time quasicrystal
  in Rydberg atomic gases,'' \emph{Nature Communications} \textbf{17},
  9186 (2026), published July 29, 2026. DOI:
  \href{https://doi.org/10.1038/s41467-026-75966-x}{10.1038/s41467-026-75966-x}.
\item
  {[}4{]} I. Martin, G. Refael and B. Halperin, ``Topological Frequency
  Conversion in Strongly Driven Quantum Systems,'' \emph{Physical Review
  X} \textbf{7}, 041008 (2017). DOI:
  \href{https://doi.org/10.1103/PhysRevX.7.041008}{10.1103/PhysRevX.7.041008}.
\item
  {[}5{]} KLOE-2 Collaboration, ``Measurement of the running of the fine
  structure constant below 1 GeV with the KLOE detector,'' \emph{Physics
  Letters B} \textbf{767}, 485--492 (2017). DOI:
  \href{https://doi.org/10.1016/j.physletb.2016.12.016}{10.1016/j.physletb.2016.12.016}.
\item
  {[}6{]} P. J. Mohr et al., ``CODATA Recommended Values of the
  Fundamental Physical Constants: 2022,'' NIST / \emph{Reviews of Modern
  Physics} (2025). \href{https://physics.nist.gov/constants}{NIST
  Constants}.
\item
  {[}7{]} A. Hurwitz, ``Über die Composition der quadratischen Formen
  von beliebig vielen Variablen,'' \emph{Nachrichten von der
  Gesellschaft der Wissenschaften zu Göttingen,
  Mathematisch-Physikalische Klasse}, 309--316 (1898).
\item
  {[}8{]} J. F. Adams, ``On the non-existence of elements of Hopf
  invariant one,'' \emph{Annals of Mathematics} \textbf{72}, 20--104
  (1960). DOI: \href{https://doi.org/10.2307/1970147}{10.2307/1970147}.
\item
  {[}9{]} D. Shechtman, I. Blech, D. Gratias and J. W. Cahn, ``Metallic
  Phase with Long-Range Orientational Order and No Translational
  Symmetry,'' \emph{Physical Review Letters} \textbf{53}, 1951--1953
  (1984). DOI:
  \href{https://doi.org/10.1103/PhysRevLett.53.1951}{10.1103/PhysRevLett.53.1951}.
\item
  {[}10{]} J. Bamberg, G. Cairns and D. Kilminster, ``The
  crystallographic restriction, permutations, and Goldbach's
  conjecture,'' \emph{The American Mathematical Monthly} \textbf{110},
  202--209 (2003). DOI:
  \href{https://doi.org/10.1080/00029890.2003.11919956}{10.1080/00029890.2003.11919956}.
\item
  {[}11{]} M. H. Jensen, P. Bak and T. Bohr, ``Transition to chaos by
  interaction of resonances in dissipative systems. I. Circle maps,''
  \emph{Physical Review A} \textbf{30}, 1960--1969 (1984). DOI:
  \href{https://doi.org/10.1103/PhysRevA.30.1960}{10.1103/PhysRevA.30.1960}.
\item
  {[}12{]} D. Gabor, ``Theory of communication. Part 1: The analysis of
  information,'' \emph{Journal of the Institution of Electrical
  Engineers -- Part III} \textbf{93}, 429--441 (1946). DOI:
  \href{https://doi.org/10.1049/ji-3-2.1946.0074}{10.1049/ji-3-2.1946.0074}.
\item
  {[}13{]} D. Gottesman, ``The Heisenberg Representation of Quantum
  Computers,'' arXiv:
  \href{https://arxiv.org/abs/quant-ph/9807006}{quant-ph/9807006}
  (1998).
\item
  {[}14{]} H. S. M. Coxeter, ``Integral Cayley numbers,'' \emph{Duke
  Mathematical Journal} \textbf{13}, 561--578 (1946). DOI:
  \href{https://doi.org/10.1215/S0012-7094-46-01347-6}{10.1215/S0012-7094-46-01347-6}.
\item
  {[}15{]} M. S. Viazovska, ``The sphere packing problem in dimension
  8,'' \emph{Annals of Mathematics} \textbf{185}, 991--1015 (2017). DOI:
  \href{https://doi.org/10.4007/annals.2017.185.3.7}{10.4007/annals.2017.185.3.7}.
\item
  {[}16{]} V. Elser and N. J. A. Sloane, ``A highly symmetric
  four-dimensional quasicrystal,'' \emph{Journal of Physics A:
  Mathematical and General} \textbf{20}, 6161--6167 (1987). DOI:
  \href{https://doi.org/10.1088/0305-4470/20/18/016}{10.1088/0305-4470/20/18/016}.
\item
  {[}17{]} A. B. Zamolodchikov, ``Integrals of motion and S-matrix of
  the (scaled) T = T\_c Ising model with magnetic field,''
  \emph{International Journal of Modern Physics A} \textbf{4},
  4235--4248 (1989). DOI:
  \href{https://doi.org/10.1142/S0217751X8900176X}{10.1142/S0217751X8900176X}.
\item
  {[}18{]} R. Coldea et al., ``Quantum Criticality in an Ising Chain:
  Experimental Evidence for Emergent E8 Symmetry,'' \emph{Science}
  \textbf{327}, 177--180 (2010). DOI:
  \href{https://doi.org/10.1126/science.1180085}{10.1126/science.1180085}.
\item
  {[}19{]} H. Zou et al., ``E8 Spectra of Quasi-One-Dimensional
  Antiferromagnet BaCo2V2O8 under Transverse Field,'' \emph{Physical
  Review Letters} \textbf{127}, 077201 (2021). DOI:
  \href{https://doi.org/10.1103/PhysRevLett.127.077201}{10.1103/PhysRevLett.127.077201}.
\item
  {[}20{]} M. B. Green and J. H. Schwarz, ``Anomaly cancellations in
  supersymmetric D = 10 gauge theory and superstring theory,''
  \emph{Physics Letters B} \textbf{149}, 117--122 (1984). DOI:
  \href{https://doi.org/10.1016/0370-2693(84)91565-X}{10.1016/0370-2693(84)91565-X}.
\item
  {[}21{]} D. J. Gross, J. A. Harvey, E. Martinec and R. Rohm,
  ``Heterotic String,'' \emph{Physical Review Letters} \textbf{54},
  502--505 (1985). DOI:
  \href{https://doi.org/10.1103/PhysRevLett.54.502}{10.1103/PhysRevLett.54.502}.
\item
  {[}22{]} T. Kugo and P. Townsend, ``Supersymmetry and the division
  algebras,'' \emph{Nuclear Physics B} \textbf{221}, 357--380 (1983).
  DOI:
  \href{https://doi.org/10.1016/0550-3213(83)90584-9}{10.1016/0550-3213(83)90584-9}.
\item
  {[}23{]} J. C. Baez and J. Huerta, ``Division algebras and
  supersymmetry I,'' in \emph{Superstrings, Geometry, Topology, and
  C}-algebras*, Proceedings of Symposia in Pure Mathematics \textbf{81},
  65--80 (2010). arXiv:
  \href{https://arxiv.org/abs/0909.0551}{0909.0551}.
\end{itemize}

\subsection{Appendix A: How to Reproduce the Numbers and Integer
Counts}\label{appendix-a-how-to-reproduce-the-numbers-and-integer-counts}

The attached \texttt{verify\_\allowbreak{}te14.py} computes, for the case of Chapter
7, the full enumeration of four-dimensional cells, the finite root, the
positive root of the geometric self-consistent value, and the three
differences. It uses the standard library only.

The attached \texttt{verify\_\allowbreak{}e8.py} confirms, for §10.3, the
construction of the 240 roots of E₈ (112 + 128), that all roots have
squared length 2 and integer inner products (an even lattice), that the
Gram matrix of the simple roots has determinant 1 (unimodular), the
dimension 248 and the Coxeter number 30, and that \(2\cos(\pi/5)\) is
the golden ratio, by exact arithmetic with fractions. For §7.6 it also
confirms that the Coxeter element satisfies the cyclotomic polynomial
\(\Phi_{30}\) (exponents 1, 7, 11, 13, 17, 19, 23, 29), the integer
relations of Eqs. (12a) and (12b), and the breakdown of the
four-dimensional 137 cells by coordinate type.

The attached \texttt{verify\_\allowbreak{}null\_\allowbreak{}from\_\allowbreak{}products.py} confirms, for
§11.1, that all matrices with integer entries preserving length (8 in
two dimensions, 48 in three dimensions) have finite order, that Eq.
(17c) gives points on the light cone identically for the Cayley--Dickson
dimensions \(k=1,2,4,8\) and fails for \(k=16\), and Euclid's formula,
by exact integer arithmetic.

The attached \texttt{verify\_\allowbreak{}phase\_\allowbreak{}lattice.py} computes, for §3.3(c),
§10.1 and §10.1.1, the exhaustive check of finite orders of integer
matrices, the orders of the \(Z_4\) and \(Z_6\) maps, and the rotation
orders allowed in dimensions 2--4 (\(\psi(m)\le d\)). It uses the
standard library only.

\begin{Shaded}
\begin{Highlighting}[]
\ExtensionTok{python}\NormalTok{ verify\_te14.py}
\ExtensionTok{python}\NormalTok{ verify\_phase\_lattice.py}
\ExtensionTok{python}\NormalTok{ verify\_e8.py}
\ExtensionTok{python}\NormalTok{ verify\_null\_from\_products.py}
\end{Highlighting}
\end{Shaded}

\subsection{Version History}\label{version-history}

\begin{itemize}
\tightlist
\item
  \textbf{v1.0 --- 2026-10-12}: First fixing of the research hypothesis
  of this dialogue. States the order of the author's questions; the
  general principle (non-commensurability, pull-in and incomplete
  locking of integer-selecting symmetries, and readout along the
  direction of evolution); the search proposal (search for pairs that
  appear theoretically unrelated yet lie close, and that pull in but
  never lock completely); the fine-structure constant as one case of it;
  footholds for time and uncertainty (integer counting of beats, the
  cost of a scale, integer phases \(Z_4\)/\(Z_6\)); the known example of
  the circle map; precedents for integer selection (the crystallographic
  restriction theorem, Hurwitz's theorem, E₈); the integer
  correspondences between the case and E₈ (near but not identical); the
  zero square sum and central projection as the ultimate symmetry, and
  that \(U^n=I\) and \(\sum x^2=0\) are consequences rather than
  assumptions (Chapter 11); underived conditions; prior studies; and
  falsification conditions. Even if later numerical experiments arise,
  the claims of this version will not be rewritten retroactively;
  differences will be recorded from v1.1 onward.
\end{itemize}

\end{document}
